\documentclass[12pt]{article}

% ------------------------
% Paquetes básicos
% ------------------------
\usepackage[spanish]{babel}
\usepackage[utf8]{inputenc}
\usepackage[T1]{fontenc}
\usepackage{amsmath, amssymb}
\usepackage{graphicx}
\usepackage{float}
\usepackage{hyperref}
\usepackage{geometry}
\usepackage{setspace}
\usepackage{caption}
\usepackage{subcaption}

\geometry{margin=1in}
\onehalfspacing

% ------------------------
% Información del documento
% ------------------------
\title{%
Simulación del Efecto Túnel Cuántico en Silicio mediante COMSOL Multiphysics\\
\large Reporte de Investigación Exploratoria
}

\author{%
Nombre del estudiante\\
Carrera\\
Universidad
}

\date{\today}

\begin{document}
\maketitle
\thispagestyle{empty}
\newpage

% ------------------------
\begin{abstract}
% Objetivo general del trabajo
% Metodología de simulación
% Alcance del estudio
% Resultados esperados
\end{abstract}

\newpage
\tableofcontents
\newpage

% ========================
\section{Introducción}
% Importancia de los semiconductores
% Relevancia del silicio
% Necesidad de modelos cuánticos
% Objetivo del reporte

% ========================
\section{Fundamentos de Mecánica Cuántica Aplicada}
\subsection{Naturaleza ondulatoria de las partículas}
\subsection{Ecuación de Schrödinger}
\subsection{Pozo de potencial unidimensional}
\subsection{Efecto túnel cuántico}

% ========================
\section{Semiconductores y Física del Silicio}
\subsection{Estructura cristalina del silicio}
\subsection{Bandas de energía}
\subsection{Portadores de carga}
\subsection{Barreras de potencial en dispositivos}

% ========================
\section{Modelado del Efecto Túnel}
\subsection{Descripción física del fenómeno}
\subsection{Aproximaciones y supuestos}
\subsection{Modelo matemático del túnel cuántico}
\subsection{Condiciones de frontera}

% ========================
\section{Simulación Numérica con COMSOL}
\subsection{Descripción del entorno COMSOL Multiphysics}
\subsection{Módulos utilizados}
\subsection{Definición de la geometría}
\subsection{Parámetros del material}
\subsection{Configuración de la malla}
\subsection{Implementación de la ecuación de Schrödinger}
\subsection{Condiciones iniciales y de frontera}
\subsection{Método numérico empleado}

% ========================
\section{Resultados de Simulación}
\subsection{Distribución de probabilidad}
\subsection{Dependencia con el ancho de la barrera}
\subsection{Influencia de la altura del potencial}
\subsection{Análisis de corriente por efecto túnel}
\subsection{Visualización de resultados}

% ========================
\section{Discusión}
\subsection{Interpretación física de los resultados}
\subsection{Limitaciones del modelo}
\subsection{Comparación con resultados teóricos}

% ========================
\section{Aplicaciones Tecnológicas}
\subsection{Dispositivos MOS}
\subsection{Diodos túnel}
\subsection{Escalamiento de transistores}
\subsection{Nanotecnología y electrónica moderna}

% ========================
\section{Conclusiones}
% Síntesis
% Aportes del trabajo
% Posibles extensiones futuras

% ========================
\section*{Agradecimientos}

% ========================
\begin{thebibliography}{99}
% Referencias
\end{thebibliography}

\end{document}
